% Part: turing-machines 
% Chapter: undecidability 
% Section: unsolvability-decision-problem

\documentclass[../../../include/open-logic-section]{subfiles}

\begin{document}

\olfileid{tur}{und}{uns} 
\olsection{判定问题的不可解性}

\begin{thm}
\ollabel{thm:decision-prob}
判定问题是不可解的：不存在这样的图灵机~$D$，使得当它以包含一阶逻辑语句~$!B$ 的磁带作为输入启动时，$D$~最终停机，并且当且仅当 $!B$ 有效时输出~$1$，否则输出~$0$。
\end{thm}

\begin{proof}
假设判定问题是可解的，即假设存在一个图灵机~$D$。那么我们就能如下解决停机问题。我们构造一个图灵机~$E$，当给定图灵机~$M_e$ 的编号~$e$ 和输入~$w$ 作为输入时，它计算对应的语句~$!T(M_e, w) \lif !E(M_e, w)$ 并停机，此时扫描磁带最左边的方格。那么，给定输入 $e$ 和 $w$，机器 $E \concat D$ 会先计算~$!T(M_e, w) \lif
!E(M_e, w)$，然后对该输入运行判定问题机器~$D$。当且仅当 $!T(M_e, w) \lif !E(M_e, w)$ 有效时，$D$ 停机并输出~$1$；否则输出~$0$。根据 \olref[ver]{lem:halt-if-valid} 和 \olref[ver]{lem:valid-if-halt}，$!T(M_e, w) \lif !E(M_e, w)$ 有效当且仅当 $M_e$ 在输入~$w$ 上停机。因此，给定输入 $e$ 和 $w$，当且仅当 $M_e$ 在输入~$w$ 上停机时，$E\concat D$ 停机并输出~$1$；否则停机并输出~$0$。换言之，$E \concat D$ 将能解决停机问题。但根据 \olref[hal]{thm:halting-problem}，这样的图灵机不可能存在。
\end{proof}

\begin{cor}\ollabel{cor:undecidable-sat}%
判定任一一阶逻辑语句是否可满足，是不可判定的。
\end{cor}

\begin{proof}
  假设存在一台图灵机~$S$ 能判定可满足性。那么我们就能如下解决（有效性）判定问题：给定输入语句~$!B$，将 $!B$ 右移一个方格。返回到第~1 个方格并写下符号~$\lnot$。

  现在运行图灵机~$S$。它最终会停机，并在磁带上输出~$1$（若 $\lnot !B$ 可满足）或~$0$（若 $\lnot !B$ 不可满足）。如果第~1 个方格上有~$\TMstroke$，就将其擦除；如果第~1 个方格是空白，就写下一个~$\TMstroke$，然后停机。

  这台图灵机总是会停机，并且当且仅当 $\lnot !B$ 不可满足时输出~$1$，否则输出~$0$。由于 $!B$ 有效当且仅当 $\lnot !B$ 不可满足，所以当且仅当 $!B$ 有效时该机器输出~$1$，否则输出~$0$，也就是说，它能解决（有效性）判定问题。
\end{proof}

\begin{explain}
因此，不存在一台图灵机能对“$!B$ 是否是一阶逻辑的有效语句？”这个问题总是给出正确的“是”或“否”答案。然而，\emph{确实存在}一台图灵机总能给出正确的“是”答案——只不过当答案是“否”时它只是不停机。这一结论源于一阶逻辑的可靠性与完备性定理，以及推导可以被能行枚举这一事实。
\end{explain}

\begin{thm}
  \ollabel{thm:valid-ce}%
  一阶逻辑语句的有效性是半可判定的：存在一台图灵机~$E$，当它以包含一阶逻辑语句~$!B$ 的磁带作为输入启动时，当且仅当 $!B$ 有效时，$E$~最终停机并输出~$1$；否则永不停机。
\end{thm}

\begin{proof}
  一阶逻辑所有可能的推导都可以被一个能行算法依次生成。机器~$E$ 逐一生成这些推导，当它找到一个能证明 $\Proves !B$ 的推导时，就停机并输出~$1$。根据可靠性定理，如果 $E$ 停机并输出~$1$，则必有~$\Entails !B$。根据完备性定理，如果 $\Entails !B$，则存在一个推导能证明~$\Proves !B$。由于 $E$ 系统地生成所有可能的推导，它终将找到一个证明~$\Proves !B$ 的推导，因此终将停机并输出~$1$。
\end{proof}

\end{document}